% GENERATED FILE. Edit canonical lessons, then run npm run generate:books.
% canonical-source-hash: fnv1a64:7f0fa755c8b6dcb8
% canonical-lessons: ES-C277-segun, ES-C277-es-decir, ES-C277-por-un-lado, ES-C277-en-resumen, ES-C277-practice

\chapter{Working From Sources --- Credit It, Restate It, Hold Both, Compress}
\label{ch:fuentes}
\hlchaptermodality{277}

\begin{chapteropening}
\textbf{By the end of this chapter:} \emph{I can attribute a claim in Spanish, restate it in my own words, hold two disagreeing sources in one structure, and compress both into a sentence.}

Complete ES-C277-practice: run all four moves on two sources that disagree, then say what the summary left out.
\end{chapteropening}

\section[según]{según — the word that says this is not mine}

\label{lesson:ES-C277-segun}



\begin{quote}
The last chapter of the book. It is about handling what other people have said, and it starts with the smallest and most important move: saying whose it is.
\end{quote}

\subsection*{What to know first}
\begin{itemize}
\raggedright
  \item en mi opinión — marking a claim as yours.
  \item dijo que — reporting what someone said.
\end{itemize}

\begin{sounds}
\begin{itemize}
\raggedright
  \item \emph{se-GÚN}, and the accent is written because the stress is on the last syllable. $\to$ reference
\end{itemize}
\end{sounds}

\begin{cousinweb}
\textbf{según} — \textquotedblleft{}according to.\textquotedblright{}

Latin \textbf{secundum}, \textquotedblleft{}following, next along\textquotedblright{} — the same word as English \textbf{second} and, through it, \textbf{sequence}. To say \emph{según} is literally to say \emph{following} — \emph{following García, the number is higher.}

You already have the opposite move: \emph{en mi opinión} marks a claim as \textbf{yours}. \emph{Según} marks it as \textbf{someone else's}, and the two together are the whole apparatus of attribution.
\end{cousinweb}

\begin{grammarlens}[title={What you've built}]
This is the smallest honest unit of scholarship, and it does two jobs at once.

It \textbf{credits} — the idea belongs to somebody. And it \textbf{distances} — you are reporting the claim, not endorsing it. \emph{Según él, es fácil} leaves entirely open whether you agree.

That second job is why \emph{según} is common in journalism and argument, not only in citation. It lets you put a claim on the table without picking it up.
\end{grammarlens}

\subsection*{Your turn}
Say these aloud:

\begin{itemize}
\raggedright
  \item \emph{según} plus a name plus a claim
  \item the same claim as your own, with \emph{en mi opinión}
  \item which of the two commits you to it
\end{itemize}

\begin{tcolorbox}[breakable,colback=teal!4,colframe=teal!35!black,title={Before you move on}]
What Latin word is \emph{según}, and which English word shares it? (\emph{Secundum} — \emph{second}.) What two jobs does it do at once? (Credits the source, and distances you from the claim.)
\end{tcolorbox}
\section[es decir]{es decir — saying it again, differently, on purpose}

\label{lesson:ES-C277-es-decir}



\begin{quote}
You have the source. Now the harder half: putting what it said into words that are yours.
\end{quote}

\subsection*{What to know first}
\begin{itemize}
\raggedright
  \item según — marking whose claim it is.
  \item es que… — a different \emph{es} phrase entirely.
\end{itemize}

\begin{sounds}
\begin{itemize}
\raggedright
  \item \emph{es de-CIR}. Note how close this is to \emph{es que} and how different its job is. $\to$ reference
\end{itemize}
\end{sounds}

\begin{cousinweb}
\textbf{es decir} — \textquotedblleft{}that is to say, in other words.\textquotedblright{}

Literally \emph{it is to say}: the infinitive \emph{decir} from Latin \textbf{dicere}, the root inside English \textbf{diction}, \textbf{dictate}, \textbf{verdict} and \textbf{contradict}.

It marks what follows as \textbf{the same content in different words} — and that is the definition of a \textbf{paraphrase}. Not a summary, which is shorter, and not a quotation, which is identical. The same, said again, deliberately.
\end{cousinweb}

\begin{grammarlens}[title={What you've built}]
Why a language needs a marker for this at all is worth a moment.

Without \emph{es decir}, a restatement looks like a second, separate claim. The listener counts two points where you made one. The marker tells them: \textbf{this is the same point, arriving again in a form you may find easier.}

That is also the honest test of whether you have understood a source. If you can only repeat its words, you have not understood it; if you can say it another way, you have. Paraphrase is comprehension made visible, which is why it is asked for in every examination ever set.
\end{grammarlens}

\subsection*{Your turn}
Say these aloud:

\begin{itemize}
\raggedright
  \item a claim, then \emph{es decir}, then the same claim differently
  \item why the marker matters (without it, it reads as a second point)
  \item the difference between a paraphrase and a summary
\end{itemize}

\begin{tcolorbox}[breakable,colback=teal!4,colframe=teal!35!black,title={Before you move on}]
What does \emph{es decir} mark? (The same content in different words.) How does a paraphrase differ from a summary? (A summary is shorter; a paraphrase is the same length, said differently.) Which English words share \emph{dicere}? (\emph{Diction}, \emph{verdict}, \emph{contradict}.)
\end{tcolorbox}
\section[por un lado… por otro]{por un lado… por otro — two sources in one sentence}

\label{lesson:ES-C277-por-un-lado}



\begin{quote}
One source you can report. Two that disagree is a different problem, and this is the frame for it.
\end{quote}

\subsection*{What to know first}
\begin{itemize}
\raggedright
  \item es decir — restating a source.
  \item sin embargo — turning against a point.
\end{itemize}

\begin{sounds}
\begin{itemize}
\raggedright
  \item \emph{por un LA-do … por O-tro}. $\to$ reference
\end{itemize}
\end{sounds}

\begin{cousinweb}
\textbf{por un lado… por otro (lado)} — \textquotedblleft{}on one hand… on the other.\textquotedblright{}

\emph{Lado} is Latin \textbf{latus}, \textquotedblleft{}side, flank\textquotedblright{} — English \textbf{lateral}, \textbf{bilateral}. The image is spatial in both languages: two sides of a thing, held apart so both can be seen.

Spanish, like English, usually drops the second \emph{lado}: \emph{por un lado… por otro}.
\end{cousinweb}

\begin{grammarlens}[title={What you've built}]
This is the frame that makes \textbf{synthesis} possible, and synthesis is not the same as reporting twice.

Reporting twice: \emph{Según García, sube. Según López, baja.} Two facts, no relation. Synthesising: \emph{Por un lado García dice que sube; por otro, López lo niega.} Now the two are \textbf{held in one structure}, and the reader can see that they conflict rather than having to notice it themselves.

The frame does the analytical work. Announcing two sides commits you to producing both, which is exactly the discipline that stops a summary from quietly becoming an argument for one of them.
\end{grammarlens}

\subsection*{Your turn}
Say these aloud:

\begin{itemize}
\raggedright
  \item two conflicting claims, reported flatly, one after the other
  \item the same two inside the \emph{por un lado… por otro} frame
  \item what the frame made visible that the flat version did not
\end{itemize}

\begin{tcolorbox}[breakable,colback=teal!4,colframe=teal!35!black,title={Before you move on}]
What Latin word is \emph{lado}? (\emph{Latus}, \textquotedblleft{}side\textquotedblright{} — English \emph{lateral}.) What does the frame do that reporting twice does not? (It holds the two in one structure, so the relation between them is visible.)
\end{tcolorbox}
\section[en resumen]{en resumen — the hardest thing to do with a text}

\label{lesson:ES-C277-en-resumen}



\begin{quote}
The last phrase in the book, and the hardest skill in it.
\end{quote}

\subsection*{What to know first}
\begin{itemize}
\raggedright
  \item por un lado… por otro — holding two sources.
  \item por lo tanto — the cohesion join.
\end{itemize}

\begin{sounds}
\begin{itemize}
\raggedright
  \item \emph{en re-su-MEN}. $\to$ reference
\end{itemize}
\end{sounds}

\begin{cousinweb}
\textbf{en resumen} — \textquotedblleft{}in short, in summary.\textquotedblright{}

\emph{Resumen} is Latin \textbf{resumere} — \emph{re-} (\textquotedblleft{}again\textquotedblright{}) plus \emph{sumere} (\textquotedblleft{}to take up\textquotedblright{}). To summarise is to \textbf{take it up again}, and English \textbf{resume} kept the taking-up while Spanish kept the shortening.

Note the false friend, because it will catch you: Spanish \emph{resumen} is a \textbf{summary}, not a curriculum vitae. The English \emph{résumé} took the same French word in a different direction.
\end{cousinweb}

\begin{grammarlens}[title={What you've built}]
Summarising is harder than paraphrasing, and the difference is \textbf{decision}.

A paraphrase keeps everything and changes the words. A summary \textbf{throws things away}, and every omission is a judgement about what mattered. That is why two honest summaries of one text can differ, and why summarising is the last skill on this ladder rather than the first.

It also closes the loop this book opened. It began by giving you a word for \emph{hello}. It ends by giving you the means to read what other people have written, decide what is load-bearing in it, and say that much in your own words — which is most of what it means to \emph{use} a language rather than to know one.
\end{grammarlens}

\subsection*{Your turn}
Say these aloud:

\begin{itemize}
\raggedright
  \item \emph{en resumen}, and compress the last thing you read into one sentence
  \item what you left out, and why
  \item the false friend to avoid (\emph{resumen} is not a CV)
\end{itemize}

\begin{tcolorbox}[breakable,colback=teal!4,colframe=teal!35!black,title={Before you move on}]
What does \emph{resumere} mean, taken apart? (To take up again.) What separates a summary from a paraphrase? (A summary throws things away, and every omission is a judgement.)
\end{tcolorbox}
\section[Practice]{Read two, say one}

\label{lesson:ES-C277-practice}



\begin{quote}
The last exercise in the book, and it uses all four moves in the order you would actually need them.
\end{quote}

\subsection*{Your turn}
Find two things that disagree — two people, two articles, two friends. Then run the four moves.

Say these aloud:

\begin{itemize}
\raggedright
  \item \emph{según} the first one, what it claims. You have credited it and not agreed with it.
  \item \emph{es decir}, and the same claim in words the first source did not use. If you cannot, you have not understood it yet.
  \item \emph{por un lado… por otro}, putting both inside one sentence so the disagreement is visible rather than merely present.
  \item \emph{en resumen}, and one sentence that covers both.
\end{itemize}

Say these aloud:

\begin{itemize}
\raggedright
  \item what you left out of the summary, and why it was safe to lose
  \item whether your summary favours one of the two. Check honestly — this is where a summary quietly becomes an argument.
\end{itemize}

\begin{tcolorbox}[breakable,colback=teal!4,colframe=teal!35!black,title={Before you move on}]
Name the four moves: credit it, restate it, hold both, compress. Which one proves you understood the source? (The paraphrase.) Which one requires you to throw something away? (The summary.)
\end{tcolorbox}
